A smoothing is applied before a histogram is translated into the likelihood function to remove the effect of statistical fluctuations. The used smoothing method is called "MaxVariation" - it is the default method in TRExFitter \href{https://gitlab.cern.ch/atlas-phys/exot/CommonSystSmoothingTool/-/blob/master/doc/algorithm_description.md}{https://gitlab.cern.ch/atlas-phys/exot/CommonSystSmoothingTool/-/blob/master/doc/algorithm\_description.md}. It is applied to all 2-point systematics, and all systematics with events stored in a separate TTree (not applied to weight systematics). In Int Note version 0.6 and earlier, a different smoothing method was used for the \ttbar systematics as baseline: TTbarResonance. It was found that the MaxVariation smoothing better captures the shape of the \ttbar systematics, so it is used as baseline since Int Note version 0.7 for all the main results, such as the fits in Sections \ref{sec:asimov}, \ref{sec:blinded}, \ref{sec:unblinded_Bonly}, \ref{sec:unblinded_SplusB}. Also the Appendix \ref{app:TTB_modeling} shows the \ttbar modelling systematics with the new baseline smoothing (MaxVariation). Older studies still use the previous baseline smoothing for \ttbar systematics (TTbarResonance).

